\documentclass[12pt]{article}

\usepackage[utf8]{inputenc}
\usepackage[T1]{fontenc}
\usepackage{amsmath, amssymb, amsfonts, mathtools}
\usepackage{geometry}
\geometry{margin=2.5cm}
\usepackage{setspace}
\onehalfspacing
\usepackage{bm}
\usepackage{hyperref}
\usepackage{booktabs, tabularx}
\usepackage{pdflscape}

\title{Austeridad - México}
\author{Ismael D. Valverde Ambriz \\ Kevin Jamel Sandoval Hernandez González}
\date{}

\begin{document}
\maketitle

\section*{Modelo Fiscal con Expectativas Racionales para México}

\paragraph{Sentido del modelo.}
El \emph{Open Lite} es una versión parsimoniosa de economía abierta, linealizada en torno al estado estacionario, que conserva los canales mínimos para analizar cómo una \textbf{consolidación fiscal} (vía reducción de gasto $g_t$ o aumento de ingresos $\tau_t$) afecta a la \textbf{prima de riesgo} $rp_t$. El mecanismo es directo: la consolidación mejora el balance primario $pb_t$, reduce la deuda $d_t$ y, con sensibilidad $k_{rp}>0$, la prima de riesgo cae. Vía la UIP, un $rp_t$ menor aprecia el tipo de cambio ($e_t$), modera inflación a través de traspaso cambiario y permite condiciones financieras más holgadas. El parámetro \textbf{$k_{rp}$} identifica \emph{cuánto} responde la prima de riesgo ante cambios en la deuda, por lo que el modelo es adecuado para cuantificar el efecto de consolidaciones sobre $rp_t$.

\medskip
\noindent\textbf{Convención:} todas las variables están en \emph{desvíos respecto al estado estacionario}. $\mathbb{E}_t[\cdot]$ denota expectativas racionales. 
\medskip

\noindent\textbf{Nota de coherencia (regla monetaria).} La regla de Taylor en este documento y en el código de simulación reacciona a la \textbf{depreciación} medida como variación del tipo de cambio, $\Delta e_t \equiv e_t - e_{t-1}$, no al nivel $e_t$.

%===============================
% BLOQUES DEL MODELO
%===============================

\subsection*{Bloque de demanda (IS)}
\begin{align}
\text{Brecha (IS):}\quad 
gap_t &= \rho_{gap}\, gap_{t-1} + a_{yg}\, g_{t-1} + a_{y\tau}\, \tau_{t-1} + \varepsilon^y_t \tag{1}
\end{align}
\noindent\emph{Descripción:} la brecha $gap_t$ (proxy del ciclo) depende de su inercia y de la política fiscal rezagada. Un ajuste fiscal contractivo hoy reduce la demanda futura vía $g_{t-1}$ y $\tau_{t-1}$.

\begin{align}
\text{PIB (proxy):}\quad 
y_t &= gap_t \tag{2}
\end{align}
\noindent\emph{Descripción:} en la aproximación lineal local, tratamos $y_t$ como proxy de la brecha (normalización útil para interpretar IRFs).

\subsection*{Bloque de precios y política monetaria}
\begin{align}
\text{Phillips NK:}\quad 
\pi_t &= \beta\, \mathbb{E}_t[\pi_{t+1}] + \kappa_\pi\, gap_t + \chi_{\Delta e}\, \Delta e_t + \varepsilon^\pi_t, 
\qquad \Delta e_t \equiv e_t - e_{t-1} \tag{3}
\end{align}
\noindent\emph{Descripción:} la inflación depende de expectativas, del ciclo y del traspaso del tipo de cambio (en variación).

\begin{align}
\text{Regla monetaria (Taylor abierta):}\quad 
i_t &= \phi_\pi\, \pi_t + \phi_y\, gap_t + \phi_{\Delta e}\, \Delta e_t + \varepsilon^i_t \tag{4}
\end{align}
\noindent\emph{Descripción:} la autoridad reacciona a inflación, brecha y depreciación. El término $\phi_{\Delta e}>0$ ayuda a asegurar unicidad (condiciones de Blanchard–Kahn) cuando $e_t$ es de salto por la UIP.

\subsection*{Bloque externo (UIP y tasas)}
\begin{align}
\text{UIP con prima:}\quad 
\mathbb{E}_t[e_{t+1}] &= e_t + i_t - i^*_t + rp_t + \varepsilon^e_t \tag{5}
\end{align}
\noindent\emph{Descripción:} el tipo de cambio (log) es variable de salto; una mayor prima $rp_t$ exige depreciación esperada.

\begin{align}
\text{Tasa externa:}\quad 
i^*_t &= \rho_{i^*}\, i^*_{t-1} + \varepsilon^{i^*}_t \tag{6}
\end{align}
\noindent\emph{Descripción:} dinámica exógena parsimoniosa de la tasa internacional en desvíos.

\subsection*{Bloque fiscal}
\begin{align}
\text{Gasto:}\quad 
g_t &= \alpha_{gg}\, g_{t-1} + \gamma_g\, gap_t + \varepsilon^g_t \tag{7}
\\
\text{Ingresos:}\quad 
\tau_t &= \alpha_{\tau\tau}\, \tau_{t-1} + \gamma_\tau\, gap_t + k_d\, d_{t-1} + \varepsilon^\tau_t \tag{8}
\\
\text{Primario:}\quad 
pb_t &= \tau_t - g_t \tag{9}
\end{align}
\noindent\emph{Descripción:} reglas fiscales simples en desvíos: inercia e impacto cíclico. La sensibilidad $k_d$ permite reacción de ingresos (o ajuste fiscal) a la deuda heredada.

\subsection*{Deuda pública (linealizada)}
\begin{align}
\text{Acumulación de deuda:}\quad 
d_t &= d_{t-1} - pb_t + \varepsilon^d_t \tag{10}
\end{align}
\noindent\emph{Descripción:} alrededor del estado estacionario neutral ($i=\pi$), la dinámica exacta se linealiza a esta forma. Un primario más alto ($pb_t\uparrow$) reduce $d_t$.

\subsection*{Prima de riesgo soberano}
\begin{align}
\text{Prima de riesgo:}\quad 
rp_t &= \rho_{rp}\, rp_{t-1} + k_{rp}\, d_{t-1} + \varepsilon^{rp}_t \tag{11}
\end{align}
\noindent\emph{Descripción:} \textbf{$k_{rp}$} es el parámetro clave: mide cómo \emph{cambios en la deuda} se trasladan a la \emph{prima de riesgo}. Una consolidación ($pb_t\uparrow \Rightarrow d_t\downarrow$) induce $rp_t\downarrow$ si $k_{rp}>0$.

%===============================
% VARIABLES, SHOCKS Y PARÁMETROS
%===============================

\subsection*{Glosario de variables y shocks}
\begin{itemize}
\item $gap_t$: brecha del producto (desvío).
\item $y_t$: proxy de PIB real (desvío), normalizado como $y_t=gap_t$.
\item $\pi_t$: inflación (desvío).
\item $i_t$: tasa de interés doméstica (desvío); $i^*_t$: tasa externa (desvío).
\item $e_t$: log del tipo de cambio (alza = depreciación); $\Delta e_t = e_t - e_{t-1}$.
\item $g_t$: gasto primario (desvío); $\tau_t$: ingresos (desvío); $pb_t=\tau_t-g_t$.
\item $d_t$: deuda pública como desvío respecto a su nivel de referencia/objetivo.
\item $rp_t$: prima de riesgo soberano (desvío).
\item Shocks: $\varepsilon^y_t$ (IS), $\varepsilon^\pi_t$ (Phillips), $\varepsilon^i_t$ (monetario), 
$\varepsilon^e_t$ (UIP), $\varepsilon^{i^*}_t$ (externo), $\varepsilon^{rp}_t$ (riesgo), 
$\varepsilon^g_t$ (gasto), $\varepsilon^\tau_t$ (ingresos), $\varepsilon^d_t$ (deuda).
\end{itemize}

\subsection*{Parámetros principales y calibración base}
\begin{table}[h!]\centering
\caption{Parámetros base}
\begin{tabular}{lcc}
\toprule
Parámetro & Símbolo & Valor \\
\midrule
Descuento Phillips & $\beta$ & 0.99 \\
Sensibilidad ciclo en precios & $\kappa_\pi$ & 0.05 \\
Traspaso cambiario & $\chi_{\Delta e}$ & 0.02 \\
Persistencia de la brecha & $\rho_{gap}$ & 0.40 \\
Canales fiscales en IS & $a_{yg},\, a_{y\tau}$ & 0.00,\;0.00 \\
Inercia gasto / ingresos & $\alpha_{gg},\,\alpha_{\tau\tau}$ & 0.20,\;0.30 \\
Estabilizadores cíclicos & $\gamma_g,\,\gamma_\tau$ & 0.15,\;0.20 \\
Regla fiscal vs deuda & $k_d$ & 0.02 \\
Regla monetaria & $\phi_\pi,\,\phi_y,\,\phi_{\Delta e}$ & 1.50,\;0.50,\;0.30 \\
Persistencia tasa externa & $\rho_{i^*}$ & 0.85 \\
Persistencia prima de riesgo & $\rho_{rp}$ & 0.80 \\
Canal deuda $\rightarrow$ riesgo & $k_{rp}$ & 0.005 \\
\bottomrule
\end{tabular}
\end{table}

\subsection*{Por qué este modelo responde la pregunta sobre $rp_t$}
La consolidación fiscal (shock $\varepsilon^g_t<0$ o $\varepsilon^\tau_t>0$) eleva $pb_t$ (ec.~(9)), reduce $d_t$ (ec.~(10)) y, con \textbf{$k_{rp}>0$}, disminuye la prima $rp_t$ (ec.~(11)). Esto afecta la UIP (ec.~(5)), moderando la depreciación esperada, y vía $\chi_{\Delta e}$ reduce presiones inflacionarias, lo que típicamente permite una respuesta monetaria menos contractiva (ec.~(4)). El parámetro \textbf{$k_{rp}$} es, por tanto, la pieza identificante del canal deuda–riesgo; su estimación o calibración cuantifica la respuesta de $rp_t$ ante una consolidación.

%===============================
% NUEVA SECCIÓN: CÓMO LEER Y CUANTIFICAR
%===============================

\section*{Cómo leer y cuantificar ``consolidación $\rightarrow$ $rp_t$''}

\paragraph{Lectura de IRFs paso a paso.}
Bajo un impulso de consolidación (por ingreso o por gasto):
\begin{enumerate}
    \item \textbf{Balance primario:} $pb_t \uparrow$ (ec.~(9)).
    \item \textbf{Deuda:} $d_t \downarrow$ en $t$ y subsecuentes (ec.~(10)).
    \item \textbf{Prima de riesgo:} $rp_t \downarrow$ vía $k_{rp}>0$ y persistencia $\rho_{rp}$ (ec.~(11)).
    \item \textbf{Tipo de cambio e inflación:} menor $rp_t$ reduce depreciación esperada (ec.~(5)), apreciando $e_t$; vía $\chi_{\Delta e}$, $\pi_t \downarrow$ (ec.~(3)).
    \item \textbf{Política monetaria:} con menor $\pi_t$ y menor $\Delta e_t$, la regla (ec.~(4)) tiende a ser menos contractiva.
\end{enumerate}

\paragraph{Multiplicador de riesgo de corto plazo.}
Para un impulso pequeño al primario, la linealización sugiere:
\[
\Delta d_t \approx -\, \Delta pb_t,
\qquad
\Delta rp_t \approx k_{rp}\,\Delta d_{t-1} \approx -\,k_{rp}\,\Delta pb_{t-1}.
\]
Con IRFs discretas, una métrica operativa a $h$ trimestres es
\[
\mathcal{M}_{rp|pb}(h) \equiv 
\frac{\text{IRF}(rp,h)}{\text{IRF}(pb,0)}.
\]
Valores negativos indican que una mejora del primario reduce la prima.

\paragraph{Efecto acumulado (área bajo la curva).}
Para cuantificar el alivio de riesgo a horizonte $H$:
\[
\text{AUC}_{rp}(H) \equiv \sum_{h=0}^{H} \text{IRF}(rp,h)\,,
\qquad
\text{AUC}_{rp|pb}(H) \equiv 
\frac{\sum_{h=0}^{H} \text{IRF}(rp,h)}{\text{IRF}(pb,0)}.
\]
La primera mide puntos porcentuales-acumulados en la prima; la segunda normaliza por la magnitud del esfuerzo fiscal inicial.

\paragraph{Semielasticidad deuda–riesgo.}
Si se observa la respuesta contemporánea de $rp$ a un cambio \emph{puro} en $d$ (manteniendo constantes otros shocks), la pendiente local es $k_{rp}$. En prácticas con IRFs de consolidación, un estimador empírico útil es:
\[
\widehat{k}_{rp}(h) \equiv
\frac{\text{IRF}(rp,h)}{\text{IRF}(d,h-1)}\quad (h\ge 1).
\]
Promediar $\widehat{k}_{rp}(h)$ para $h=1,\dots,H$ estabiliza el estimador cuando hay realimentaciones.

\paragraph{Comparación por instrumento fiscal.}
Para distinguir consolidación por \emph{gasto} vs \emph{ingresos}, compare:
\begin{itemize}
    \item Pico y AUC de $pb$, $d$ y $rp$.
    \item Respuesta de $gap$/$y$ (coste cíclico) y de $\pi$.
    \item Movimientos en $e$ y $\Delta e$ (canal externo).
\end{itemize}
El diseño preferible es el que produce mayor $\lvert \text{AUC}_{rp|pb}\rvert$ con menor caída en $gap$/$y$.

\paragraph{Nota práctica para reproducibilidad.}
Al correr \texttt{Dynare}, exporte IRFs y compute \(\mathcal{M}_{rp|pb}(h)\) y \(\text{AUC}_{rp|pb}(H)\) en una tabla. Esto permite reportar, por ejemplo, el \emph{riesgo-multiplicador fiscal} a 4 y 8 trimestres.

%===============================
% FIN
%===============================

\end{document}